\documentclass{report}
\usepackage{amsmath,amssymb}
\setlength{\parindent}{0mm}
\setlength{\parskip}{1em}
\begin{document}
\begin{center}
\rule{6in}{1pt} \
{\large Simon Baumann \\
{\bf Aspects of Multigrid for Mantle Convection}}

Geophysics - Department of Earth and \\ Environmental Sciences \\ Universit�t M�nchen (LMU) \\ Theresienstr 41 \\ 80333 M�nchen \\ Germany
\\
{\tt baumann@geophysik.uni-muenchen.de}\\
Marcus Mohr\end{center}

Convection in the Earth's mantle can be described by a Stokes-type equation
with strongly varying viscosity values coupled to an energy equation. The
dynamics of the mantle depend essentially on the underlying viscosity structure.
Computing the quasi-stationary flow field in each time step dominates the
computational cost of Earth mantle simulators. Therefore constructing efficient
solvers is crucial for simulating Earth mantle dynamics.
In this article we consider a geometric multigrid method for the viscous
operator of the Stokes-type system and study the convergence behaviour for
different smoothers and transfer operators. We focus on two aspects, the
influence of operator complexity, starting with a simplified form and ending
with its most general formulation, and the effects of highly discontinuous
viscosity parameters. Systematic numerical tests help to identify the most
efficient multigrid components.


\end{document}
